\section*{Chapitre \ref{ch:g4:symmetry} --- La symétrie}

\begin{solution}{exo:g4:symmetry:1}
Carré~: oui ($4$ axes). Rectangle~: oui ($2$). Triangle scalène~:
aucun pli ne fonctionne. Cercle~: oui --- toute droite passant par
le centre.
\end{solution}

\begin{solution}{exo:g4:symmetry:2}
Carré~: $4$ (deux médianes, deux diagonales). Rectangle~: $2$ (les
médianes seulement). Triangle équilatéral~: $3$ (une par chaque
sommet).
\end{solution}

\begin{solution}{exo:g4:symmetry:3}
En pliant le long de la diagonale, les deux moitiés triangulaires
ont la même forme mais ne se superposent pas --- l'une est
«~tournée du mauvais côté~»~: le calque se pose à côté de la
figure, pas dessus. Donc la diagonale n'est pas un axe.
\end{solution}

\begin{solution}{exo:g4:symmetry:4}
Le L image est à $2$ carreaux de \emph{l'autre} côté de l'axe,
écrit à l'envers (un L miroir), mêmes dimensions.
\end{solution}

\begin{solution}{exo:g4:symmetry:5}
La figure complétée montre le drapeau et son reflet à l'envers sous
l'axe, mât sous mât.
\end{solution}

\begin{solution}{exo:g4:symmetry:6}
$P$ ne bouge pas~: plier le long de l'axe laisse tout point
\emph{sur} l'axe exactement où il est --- $P$ est sa propre image.
\end{solution}

\begin{solution}{exo:g4:symmetry:7}
Axe vertical~: A, H, O, T. Axe horizontal~: B, E, H, O. Les deux~:
H, O. Aucun~: F, N.
\end{solution}

\begin{solution}{exo:g4:symmetry:8}
Oui~: le triangle est isocèle, et son axe passe par le sommet entre
les deux côtés de $5$~cm et par le milieu de la base de $3$~cm. Le
pli échange les deux côtés égaux.
\end{solution}

\begin{solution}{exo:g4:symmetry:9}
Toute figure avec exactement deux axes convient~: un rectangle (non
carré), un signe plus aux bras de deux longueurs différentes
(horizontal long, vertical court), un ovale. Les deux axes se
croisent toujours au centre de la figure.
\end{solution}

\begin{solution}{exo:g4:symmetry:10}
La maison complète est la moitié gauche plus son jumeau miroir~: un
mur deux fois plus large que le demi-mur, sous un toit triangulaire
complet. Chaque longueur à droite égale son partenaire à gauche ---
la symétrie copie exactement les longueurs.
\end{solution}

\begin{solution}{exo:g4:symmetry:11}
Les dessins soutiennent les deux sens~: le triangle isocèle se plie
le long de la droite passant par son sommet (deux côtés égaux
$\to$ axe)~; l'équilatéral se plie de trois façons~; le scalène n'a
pas d'axe (pas de côtés égaux $\to$ pas d'axe). La double
affirmation d'Aya est juste --- elle sera prouvée dans
\cref{ch:g6:symmetry}.
\end{solution}
